\chapter{RESULTS AND DISCUSSION} \label{c6}
\section{Testing Strategy}
Testing was conducted on multiple Android devices for teacher and student workflows. The key scenarios validated include session creation, OTP mismatch, session expiry, Bluetooth failure, face mismatch, successful attendance submission, and duplicate attendance prevention.

\section{Observed Results}
\subsection{Teacher Module}
Teacher module successfully generated live OTP, displayed timer, and managed session lifecycle. Live student list under ongoing session was refreshed with marked attendance records.

\subsection{Student Module}
Student module correctly enforced verification sequence. Attendance submission was accepted only after valid OTP, Bluetooth success, and face verification success.

\subsection{Backend Persistence}
Attendance records were inserted with required fields: student identification, session mapping, class date and period, and timestamp. Duplicate marking attempts were blocked.

\section{Result Summary Table}
\begin{table}[H]
\centering
\caption{Functional Test Outcomes}
\label{c6:tab1}
\begin{tabular}{|L{4.4cm}|L{2.4cm}|L{3.8cm}|}
\hline
\textbf{Test Case} & \textbf{Expected} & \textbf{Outcome}\\
\hline
Teacher starts session & OTP + timer visible & Passed\\
\hline
Student enters invalid OTP & Error message & Passed\\
\hline
Expired session submission & Reject attendance & Passed\\
\hline
Bluetooth verification failure & Block marking & Passed\\
\hline
Face verification failure & Block marking & Passed\\
\hline
All checks success & Mark attendance & Passed\\
\hline
Duplicate marking same slot & Reject duplicate & Passed\\
\hline
\end{tabular}
\end{table}

\section{Performance and Usability Discussion}
The guided workflow reduced user confusion and made attendance marking predictable. Teacher-side one-tap session control reduced unnecessary UI complexity. Student-side staged verification reduced invalid submissions and improved transparency in failure reasons.

\section{Limitations}
The system still depends on camera quality, lighting conditions, and Bluetooth signal behavior across devices. Network interruptions can delay synchronization in some conditions. These limitations are manageable with better fallback and buffering mechanisms in future versions.

\section{Practical Impact}
The system demonstrates a practical, deployable classroom attendance mechanism that balances security and usability. It minimizes manual effort and significantly improves attendance credibility.

\section{Conclusion}
Testing and results confirm that the implemented multi-factor attendance workflow is stable, effective, and aligned with the project objectives.

